\subsection{拉格朗日 $A$-膜与量子化}\label{abq}

\subsubsection{Hilbert 空间的定义}\label{hilbdef}

现在引入一个拉格朗日 $A$-膜，并解释它与量子化的关系。我们沿用 \cite{GW} 第 2.3 节的说明。一些重要的示意性例子此前已由 \cite{AZ} 分析过。

要量子化实辛流形 $M$，选取支撑为 $M$、带某个 $\CP$ 丛的秩 1 拉格朗日 $A$-膜 $\B$。秩 1 拉格朗日 $A$-膜的特征是其 $\CP$ 丛的曲率 $\sF$ 满足 $\sF+\sB=0$。由于已取 $\sB=\omega_J$，曲率必须是 $\sF=-\omega_J$。如第 \ref{amo} 节所述，我们的目标是用辛结构 $\omega_J$ 量子化 $M$。回忆几何量子化的第一步是引入曲率等于辛形式（此处为 $\omega_J$）的预量子线丛 $\frL$。由于 $\sF$ 是 $\frL$ 曲率的负值，可以把 $\B$ 的 $\CP$ 丛解释为 $\frL^{-1}$。这个等同合理的一个迹象是反常匹配：几何量子化中松散地称为预量子线丛的东西，更准确地说是 $\spinc$ 结构（第 \ref{gq} 节讨论过）；类似地，秩 1 膜的 $\CP$ 丛其实也是 $\spinc$ 结构而非复线丛，如脚注 \ref{anomdesc} 所述。

于是，带预量子线丛 $\frL$ 的 $M$ 的膜量子化的 Hilbert 空间定义为
\[
\H=\Hom(\B,\B_\cc),
\]
即 $(\B,\B_\cc)$ 系统的物理态空间。对偶空间是 $\H'=\Hom(\B_\cc,\B)$。把 $\H$ 而非 $\H'$ 认作 Hilbert 空间，是因为 $\A=\Hom(\B_\cc,\B_\cc)$ 从左边作用在 $\H=\Hom(\B,\B_\cc)$ 上，而从右边作用在 $\H'=\Hom(\B_\cc,\B)$ 上。若可能，人们总倾向于用环的左作用，因此如上定义 $\H$ 是自然的。把 $\H$ 称为 Hilbert 空间而不只是向量空间，是指 $\H$ 有自然的 Hermite 内积；这将在第 \ref{hermitian} 节描述。

为理解这个量子 Hilbert 空间定义的动机，我们构造开弦低能态的有效作用量，该弦的左、右边界分别落在 $\B_\cc$ 与 $\B$ 上。设弦的世界面为 $\R\times I$，$\R$ 以 $\tau$ 参数化，$I$ 为区间。$I$ 两端的边界条件使得没有费米零模；此陈述的解释见附录 \ref{details}。玻色零模沿 $I$ 为常数，故描述映射 $X:\R\to M$，可用函数 $x^i(\tau)$ 表示，$x^i$ 为 $M$ 的局部坐标。$x^i(\tau)$ 的有效作用量是什么？体作用量给出若干关于 $\d x^i/\d\tau$ 二次的项，但最相关的项关于 $\d x^i/\d\tau$ 线性，来自弦两端与 $\CP$ 丛的边界耦合。回顾式 \eqref{kogo}，在 $\Hom(\B,\B_\cc)$ 的定义中，弦的一端是 $\B_\cc$ 的 $\CP$ 丛，另一端是 $\B$ 的 $\CP$ 丛的对偶（即逆）。$\B_\cc$ 的 $\CP$ 丛完全平凡，$\B$ 的 $\CP$ 丛为 $\frL^{-1}$，其对偶为预量子线丛 $\frL$。若 $\sA$ 是 $\frL$ 上的联络，则弦端点与 $\frL$ 的耦合对低能有效作用量的贡献就是
\be\label{kixx}\sI=\int_\R x^*(\sA).\ee
这正是第 \ref{uses} 节引入的作用量，那里预量子线丛上的联络记作 $\sT/\hbar$。如那里所解释的，形式上这是与量子化相联系的作用量。

如第 \ref{uses} 节所讨论的，作用量 \eqref{iff} 或 \eqref{kixx} 没有更多信息其实无法量子化。正确的说法不是 $\Hom(\B,\B_\cc)$ 的 $A$-模型计算归结为作用量 \eqref{kixx} 的量子化；而是：$A$-模型（对适当的 $Y$）良定义而式 \eqref{kixx} 的量子化有歧义，这意味着在约化到 \eqref{kixx} 时丢失了一些信息。结论的正确表述是：$\H=\Hom(\B,\B_\cc)$ 就是 $A$-模型给出的、带预量子线丛 $\frL$ 的 $M$ 的量子化答案。

\subsubsection{若干问题}\label{ques}

说 $A$-模型给出了 $M$ 量子化的这个答案，会引出若干问题。以下是一些明显的问题：

(1) $A$-模型的物理态空间通常不是 Hilbert 空间。它通常没有任何类型的内积（甚至不定内积也没有）。在目前语境中，$A$-模型如何产生 Hermite 内积？

(2) 作用在量子 Hilbert 空间 $\H$ 上的可观测量是什么？

(3) 当标准量子化方法可用时，这个概念与它们相容吗？

第一个问题将在第 \ref{hermitian} 节回答。

关于第二个问题，$A$-模型直接给出的答案是：$\A=\Hom(\B_\cc,\B_\cc)$ 作用在 $\H=\Hom(\B,\B_\cc)$ 上。若 $M$ 是实仿射簇、$Y$ 是其复化，这个答案最令人满意。此时如第 \ref{abranes} 节所述，$\A$ 是 $Y$ 上多项式函数环 $\A_0$ 的结合非交换形变。$Y$ 上的多项式函数环与 $M$ 上的多项式函数环相同（$M$ 上的多项式函数可解析延拓为 $Y$ 上的多项式函数；$Y$ 上的非零多项式限制到 $M$ 上非零）。所以被量子化的 $M$ 上函数是多项式函数。这是相当令人满意的答案。例如若 $M=\R^{2n}$，被量子化的函数就是坐标与动量的多项式。在第 \ref{hermitian} 节我们将看到，$\A$ 中对应 $\H$ 上 Hermite 算符的元素正是 $\tau$-不变元。

然而，用膜量子化辛流形 $M$ 对实仿射簇之外的许多例子也有意义。在许多其它情形中 $\A_0$ 小得多，于是得不到如此丰富的、可量子化为 $\H$ 上算符的 $M$ 上函数。此时一般而言，什么是有趣的可观测量就不那么明显了。不过还有另一种定义作用在 $\H$ 上算符的办法，即与拉格朗日对应相联系的算符。这将在第 \ref{symq} 节解释。例如，在与几何 Langlands 相关、将在别处研究的情形 \cite{GaW} 中，$\A_0$ “很小”，$\A$ 实际上是交换的且同构于 $\A_0$。在那个特例中，重要的是考虑与拉格朗日对应相联系的算符（确切地说是 Hecke 对应，物理上与 't Hooft 算符相联系）。

最后一个问题——把膜量子化与标准量子化方法作比较——是第 \ref{comparing} 节的主题。

\subsubsection{推广}\label{genera}

到目前为止，我们考虑的拉格朗日膜的支撑为 $M$，即反全纯对合 $\tau:Y\to Y$ 的不动点集的一个分支。如第 \ref{hermitian} 节将讨论的，这个条件对在空间 $\H$ 上定义 Hermite 内积很重要。但其它类型的拉格朗日膜也重要。

由于把 $Y$ 视为辛结构为 $\omega_K$ 的辛流形，按定义，说 $L$ 是拉格朗日子流形意味着 $\omega_K$ 限制到 $L$ 上为零。但 $\omega_J$ 限制到 $L$ 上没有一般性约束。若该限制也为零，则 $L$ 实际上是关于 $Y$ 的复辛结构 $\Omega_I=\omega_J+\i\omega_K$ 的拉格朗日子流形。这样的 $L$ 是 $Y$ 的复子流形\footnote{若 $\omega_J$ 与 $\omega_K$ 限制到 $L$ 上均为零，这意味着沿 $L$，它们都是 $L$ 的切丛与法丛之间的配对。由此可知 $I=\omega_K^{-1}\omega_J$ 把 $L$ 的切丛映到自身，故 $L$ 是复子流形。}，所以它其实是复拉格朗日子流形。由于对任何 $A$-膜 $\B$，$\Hom(\B,\B_\cc)$ 都是 $\A=\Hom(\B_\cc,\B_\cc)$ 的模，特别地当 $\B$ 是支撑在复辛流形 $L$ 上的拉格朗日膜时也是如此。这类膜是 $(B,A,A)$ 型的：它们关于复结构 $I$ 是 $B$-膜（因为 $L$ 是复子流形，而 $\B$ 的 $\CP$ 丛平坦，故关于 $I$ 全纯），关于辛结构 $\omega_J$ 与 $\omega_K$ 是 $A$-膜（因为 $L$ 关于这两个结构都是拉格朗日的，且其 $\CP$ 丛平坦）。与 $(B,A,A)$ 型膜相联系的模很有意思，并将在第 \ref{comparing} 节成为重要工具，但它们与 $L$ 的量子化没有密切关系。确实，由于 $\omega_J$ 与 $\omega_K$ 限制到 $L$ 上都为零，上述讨论没有在 $L$ 上给出辛结构，故我们没有量子化 $L$ 的自然问题。（若 $Y$ 是超 Kähler 的，它还有第三个辛结构 $\omega_I$，而复拉格朗日子流形 $L$ 关于 $\omega_I$ 是辛的。但我们研究的 $A$-模型以及膜 $\B_\cc$、$\B$ 对 $\omega_I$ 不敏感。）

与 $L$ 的量子化相关的是相反情形：$\omega_J$ 限制到 $L$ 上非退化。若 $\B$ 是支撑在任一这样的 $L$ 上的秩 1 $A$-膜，则 $\Hom(\B,\B_\cc)$ 具有人们期待的、关于辛结构 $\omega=\omega_J|_L$ 的 $L$ 的量子化的大部分性质。构造中只有一步失败，即 Hermite 内积的定义。如第 \ref{hermitian} 节将解释的，这个内积的定义要求 $L$ 是 $Y$ 的某个反全纯对合的不动点集。直观上应当如此：若 $L$ 不是 $Y$ 的反全纯对称的不动点集，则 $Y$ 上的全纯函数限制到 $L$ 上何时为实并不自然，于是没有明显办法判定 $Y$ 上哪些全纯函数应在 $L$ 的量子化中成为 Hermite 算符。因此应当不能在 $\Hom(\B,\B_\cc)$ 上定义 Hermite 内积。

在 \cite{GW} 第 3.7 节中描述了 $\SL(2,\C)$ 的余伴随轨道（即式 \eqref{cosu} 定义的仿射簇 $Y$）的拉格朗日子流形 $L$ 的例子，使 $\omega_J$ 限制到 $L$ 上非退化，但 $L$ 不是 $Y$ 的任何反全纯对合的不动点集。这样的 $L$ 例如与 $\SL(2,\R)$ 或其 Lie 代数的各种有趣但非酉的表示相联系。

另一个推广是考虑 $\B$ 为支撑在 $M$（或更一般的拉格朗日子流形 $L$）上、秩 $m>1$ 的拉格朗日膜。讨论此情形时，为简单起见假定 $M$ 是 spin 流形，从而不必讨论通常涉及 spin 的反常。于是 $M$ 的 $\CP$ 丛是秩 $m$ 的平坦向量丛 $\U\to M$，结构群为 $\UU(m)$。在此情形下，$\Hom(\B,\B_\cc)$ 将是量子化概念的一种秩 $m$ 推广。这种量子化推广的存在或许稍令人惊讶。不过不难看出，在几何量子化可用的情形中这样的推广确实存在。例如设 $M=T^*N$。则按几何量子化的精神，可考虑由 $\K^{1/2}\otimes\U\to N$ 的 $\ltwo$ 截面构成的 Hilbert 空间。或者设 $M$ 有 Kähler 极化，则可考虑 Hilbert 空间 $H^0(M,K^{1/2}\otimes\U)$。
